\chapter{Mathematical background}
\label{ch:background}

\section{The splitting framework}

Let $A$ be nonsingular and write $A = M - N$ with $M$ invertible and cheap to
invert. Then $Ax = b$ is equivalent to $Mx = Nx + b$, and the associated
iteration is

\begin{equation}
  x_{k+1} = M^{-1}\left(N x_k + b\right) = x_k + M^{-1} r_k,
  \qquad r_k = b - A x_k .
  \label{eq:splitting}
\end{equation}

Subtracting the exact solution gives the error recurrence
$e_{k+1} = G e_k$ with iteration matrix

\begin{equation}
  G = M^{-1} N = I - M^{-1} A .
  \label{eq:iteration-matrix}
\end{equation}

The iteration converges for every starting vector if and only if
$\rho(G) < 1$, and the asymptotic rate is $\rho(G)$: the error is reduced by
that factor per iteration once the dominant mode dominates
\cite[Ch.~4]{saad2003}. Everything in this chapter, apart from conjugate
gradient, is a choice of $M$.

Throughout, $A = D + L + U$ splits $A$ into its diagonal, strictly lower and
strictly upper triangular parts.

\section{The model problem}

The test operator is the five point discretisation of $-\Delta u = f$ on the
unit square with homogeneous Dirichlet data. On a uniform mesh with spacing
$h = 1/(n+1)$, multiplying through by $h^2$ gives, for interior indices,

\begin{equation}
  4 u_{i,j} - u_{i-1,j} - u_{i+1,j} - u_{i,j-1} - u_{i,j+1} = h^2 f_{i,j} .
  \label{eq:stencil}
\end{equation}

The resulting matrix is symmetric positive definite with eigenvalues
$4 - 2\cos(p \pi h) - 2\cos(q \pi h)$ for $p, q = 1, \dots, n$, and
eigenvectors $\sin(p \pi x)\sin(q \pi y)$. Its extreme eigenvalues are
$O(h^2)$ and close to $8$, so the spectral condition number is $O(h^{-2})$.
This is the canonical model problem precisely because these quantities are known
in closed form, which lets every convergence claim below be checked rather than
believed \cite[Sec.~11.2]{golub2013}.

\paragraph{A trap worth recording.} The obvious manufactured solution
$u = \sin(\pi x)\sin(\pi y)$ is \emph{exactly} the lowest eigenvector of
\eqref{eq:stencil}. With a zero initial guess the residual is then a single
eigenvector, the Krylov subspace conjugate gradient builds is one dimensional,
and the method converges in one iteration at any grid size. That measures the
degeneracy of the problem, not the speed of the method. The same choice is ideal
for the stationary methods, because the slowest decaying mode of the Jacobi
iteration matrix is that same lowest mode. Both right hand sides are therefore
provided and each is used where it is honest: the single mode source for
discretisation error and stationary rates, a seeded spectrally rich source for
anything involving conjugate gradient and for the whole benchmark sweep.

\section{Richardson iteration}

Take $M = I/\omega$. Then \eqref{eq:splitting} collapses to

\begin{equation}
  x_{k+1} = x_k + \omega\left(b - A x_k\right),
\end{equation}

with $G = I - \omega A$. For symmetric positive definite $A$ with eigenvalues in
$[\lambda_{\min}, \lambda_{\max}]$ the eigenvalues of $G$ are
$1 - \omega\lambda$, so the method converges exactly when
$0 < \omega < 2/\lambda_{\max}$, and $\rho(G)$ is minimised at
$\omega = 2/(\lambda_{\min} + \lambda_{\max})$, where it equals
$(\kappa - 1)/(\kappa + 1)$ \cite[Sec.~4.1]{saad2003}.

The library chooses $\omega = 1/\|A\|_{\text{Gershgorin}}$ when none is given.
Gershgorin's bound overestimates $\lambda_{\max}$ by construction, so the
resulting step is always strictly inside the convergence interval. An earlier
implementation estimated $\lambda_{\max}$ by power iteration and diverged,
because the Rayleigh quotient approaches $\lambda_{\max}$ from below and an
unconverged estimate therefore yields a step that is too large.

\section{Jacobi}

Take $M = D$, so $N = -(L+U)$ and componentwise

\begin{equation}
  x_i^{(k+1)} = \frac{1}{a_{ii}}\left(b_i - \sum_{j \neq i} a_{ij} x_j^{(k)}\right).
\end{equation}

Every component of the new iterate depends only on the old one, which is what
makes the sweep embarrassingly parallel and why its result cannot depend on how
the unknowns were partitioned.

On the model problem the Jacobi iteration matrix has spectral radius

\begin{equation}
  \rho_J = \cos(\pi h),
  \label{eq:rho-jacobi}
\end{equation}

so $\rho_J = 1 - O(h^2)$ and the iteration count to a fixed tolerance grows like
$h^{-2}$, that is, like the number of unknowns. Convergence is guaranteed for
any strictly diagonally dominant matrix \cite[Thm.~4.9]{saad2003}.

\section{Gauss Seidel, and the factor of two}

Take $M = D + L$. Then $(D+L)x_{k+1} = b - Ux_k$ is solved by forward
substitution:

\begin{equation}
  x_i^{(k+1)} = \frac{1}{a_{ii}}\left(b_i
    - \sum_{j<i} a_{ij} x_j^{(k+1)}
    - \sum_{j>i} a_{ij} x_j^{(k)}\right).
\end{equation}

The new iterate is used as soon as it is available, which is why the method
converges faster and why the sweep is sequentially dependent.

For a consistently ordered matrix with property A, of which the five point
stencil in natural ordering is the archetype, Young's theorem gives

\begin{equation}
  \rho_{GS} = \rho_J^{\,2}
  \label{eq:rho-gs}
\end{equation}

\cite[Thm.~4.3]{young1971}. Since the iteration count to a fixed tolerance is
$\log(\varepsilon)/\log\rho$, this means Gauss Seidel needs asymptotically
\emph{half} as many iterations as Jacobi. Section \ref{sec:convergence-results}
reports the measured ratio.

Taking $M = D + U$ instead gives the backward sweep. On a symmetric matrix the
two iteration matrices are similar, so they converge at identical rates; they
differ in how information propagates across the grid, which is why the composite
of one forward and one backward sweep is worth having. That composite,
symmetric Gauss Seidel, has preconditioner
$M = (D+L)D^{-1}(D+U)$, which is symmetric whenever $A$ is. That symmetry is the
whole point: it is what a conjugate gradient preconditioner requires and what
the unsymmetric single sweep cannot provide.

\section{Successive over relaxation}

Take $M = D/\omega + L$, giving the Gauss Seidel update extrapolated by $\omega$:

\begin{equation}
  x_i^{(k+1)} = (1-\omega) x_i^{(k)}
    + \frac{\omega}{a_{ii}}\left(b_i
      - \sum_{j<i} a_{ij} x_j^{(k+1)}
      - \sum_{j>i} a_{ij} x_j^{(k)}\right).
\end{equation}

Kahan's theorem gives the necessary condition $0 < \omega < 2$, and the theorem
of Ostrowski and Reich makes it sufficient for symmetric positive definite $A$
\cite[Sec.~4.2]{saad2003}. For a consistently ordered matrix Young's theory
gives the optimum in closed form:

\begin{equation}
  \omega^{*} = \frac{2}{1 + \sqrt{1 - \rho_J^{\,2}}},
  \qquad
  \rho(G_{\omega^{*}}) = \omega^{*} - 1 .
  \label{eq:omega-opt}
\end{equation}

On the model problem $\rho_J = \cos(\pi h)$, so
$\omega^{*} = 2/(1 + \sin(\pi h))$ and the asymptotic rate becomes
$1 - O(h)$ instead of $1 - O(h^2)$. This is a change in the \emph{order} of the
iteration count, from $O(n^2)$ to $O(n)$ on an $n$ by $n$ grid, and it is the
single largest improvement available inside the stationary family
\cite[Ch.~4 to 6]{young1971}.

\section{Red black ordering}

The graph of the five point stencil is bipartite: colour cell $(i,j)$ red when
$i+j$ is even and black when it is odd, and every red cell has only black
neighbours. Reordering by colour puts $A$ in the form

\begin{equation}
  \begin{pmatrix} D_r & C \\ C^{T} & D_b \end{pmatrix},
\end{equation}

with both diagonal blocks diagonal. A sweep over all red cells is therefore a
Jacobi update within that colour and is fully parallel, and the red and black
half sweeps together are exactly Gauss Seidel in the red black ordering
\cite[Sec.~12.4]{saad2003}. This is the only Gauss Seidel style method a GPU can
run, since the natural ordering recurrence offers no parallelism at all
\cite[Ch.~6]{hager2010}.

The reordering changes the iteration matrix, so it converges at a slightly
different, generally slightly worse, per iteration rate. That penalty is
measured in Section \ref{sec:convergence-results} rather than asserted. The red
black ordering remains consistently ordered in Young's sense, so
\eqref{eq:omega-opt} carries over unchanged, which is why red black SOR is the
standard choice for parallel stencil codes.

\section{Block methods}

Partition the unknowns into contiguous blocks and write $A$ in the corresponding
block form. Taking $M$ to be the block diagonal of $A$ gives block Jacobi, in
which each block solves

\begin{equation}
  A_{pp}\, x_p^{(k+1)} = b_p - \sum_{q \neq p} A_{pq}\, x_q^{(k)}
\end{equation}

independently of the others; taking $M$ to be the block lower triangle gives
block Gauss Seidel, in which blocks are visited in order and read the already
updated ones. Point methods are the special case of one unknown per block.

Each diagonal block is solved \emph{exactly}, by a direct method matched to its
structure: dense LU with partial pivoting for the dense problem, and the Thomas
algorithm for the model problem, whose natural block is a grid line and whose
diagonal block is therefore tridiagonal and diagonally dominant. On the model
problem this is line relaxation, and solving each line exactly rather than each
point cuts the iteration count by roughly a factor of two
\cite[Sec.~4.1.1]{saad2003}.

The block count is a solver parameter and never the worker count. That is what
keeps the iterates of the block methods independent of how many workers are
running, which in turn is what lets the equivalence suite compare them bit for
bit across backends.

\section{Conjugate gradient}

For symmetric positive definite $A$, solving $Ax = b$ is equivalent to
minimising $f(x) = \tfrac{1}{2}x^{T}Ax - b^{T}x$, whose gradient is $Ax - b$, so
the residual is the direction of steepest descent. Steepest descent is slow
because successive steps undo one another; conjugate gradient fixes this by
choosing search directions that are $A$ orthogonal, $p_i^{T} A p_j = 0$ for
$i \neq j$. Under that condition an exact line search along each direction never
needs revisiting, so after $k$ steps the iterate is optimal over the Krylov
subspace $\mathrm{span}\{r_0, Ar_0, \dots, A^{k-1}r_0\}$ and, in exact
arithmetic, the method terminates in at most $n$ steps \cite{shewchuk1994}.

The practical recurrence needs one matrix vector product and two inner products
per iteration:

\begin{align}
  \alpha_k   &= \frac{r_k^{T} r_k}{p_k^{T} A p_k}, &
  x_{k+1}    &= x_k + \alpha_k p_k, &
  r_{k+1}    &= r_k - \alpha_k A p_k, \notag \\
  \beta_k    &= \frac{r_{k+1}^{T} r_{k+1}}{r_k^{T} r_k}, &
  p_{k+1}    &= r_{k+1} + \beta_k p_k. &&
  \label{eq:cg}
\end{align}

The error in the $A$ norm satisfies

\begin{equation}
  \|e_k\|_A \le 2\left(\frac{\sqrt{\kappa}-1}{\sqrt{\kappa}+1}\right)^{k}\|e_0\|_A .
  \label{eq:cg-bound}
\end{equation}

On the model problem $\kappa = O(h^{-2})$, so $\sqrt{\kappa} = O(h^{-1})$ and the
iteration count is $O(n)$: the same order as optimally relaxed SOR, but reached
without needing to know a relaxation factor in advance \cite[Ch.~6]{saad2003}.
That is why it displaced the stationary methods in practice.

The two inner products are global reductions, and they are the reason conjugate
gradient scales worse than Jacobi on a distributed machine: each one is a
synchronisation point that no amount of local work can hide. Section
\ref{sec:mpi-results} measures that cost.

\section{The numerics module}

The solvers rest on a small set of classical routines, each with its
convergence order documented and tested: bisection, Newton and Brent for root
finding \cite[Ch.~2]{burden2016}\cite[Ch.~4]{brent1973}; adaptive Simpson,
Gauss Legendre with nodes computed by Newton iteration on the Legendre
recurrence rather than tabulated, and Romberg extrapolation
\cite[Ch.~4]{burden2016}\cite{davis1984}; classical RK4 and adaptive Dormand
Prince 5(4) \cite{dormand1980}\cite[Sec.~II.4]{hairer1993}; and dense LU with
partial pivoting together with Householder QR
\cite[Sec.~3.2, 5.2]{golub2013}\cite[Lec.~10, 16]{trefethen1997}.

Dormand Prince is preferred to Fehlberg because its coefficients are fitted to
minimise the error of the fifth order formula, so the fifth order result is the
one propagated and the embedded fourth order formula serves only as the error
estimate; Fehlberg optimises the fourth order formula instead and thereby wastes
the better solution. Householder reflections are preferred to Gram Schmidt
because they are unconditionally backward stable, whereas classical Gram Schmidt
loses orthogonality in proportion to the condition number.
